\documentclass[11pt]{article}
\usepackage[T1]{fontenc}
\usepackage[margin=1in]{geometry}
\usepackage{amsmath,amssymb,amsthm,mathtools,hyperref}
\usepackage{mathrsfs}
\hypersetup{hidelinks}
\newtheorem{theorem}{Theorem}[section]
\newtheorem{lemma}[theorem]{Lemma}
\newtheorem{proposition}[theorem]{Proposition}
\newtheorem{corollary}[theorem]{Corollary}
\theoremstyle{definition}
\newtheorem{remark}[theorem]{Remark}
\newcommand{\Q}{\mathbf Q}
\newcommand{\C}{\mathbf C}
\newcommand{\Z}{\mathbf Z}
\newcommand{\F}{\mathbf F}
\newcommand{\SL}{\operatorname{SL}}
\newcommand{\GL}{\operatorname{GL}}
\newcommand{\Sp}{\operatorname{Sp}}
\newcommand{\PGL}{\operatorname{PGL}}
\newcommand{\Hg}{\operatorname{Hg}}
\newcommand{\MT}{\operatorname{MT}}
\newcommand{\Hdg}{\operatorname{Hdg}}
\newcommand{\Div}{\operatorname{Div}}
\newcommand{\End}{\operatorname{End}}
\newcommand{\Hom}{\operatorname{Hom}}
\newcommand{\Gal}{\operatorname{Gal}}
\newcommand{\Cent}{\operatorname{Cent}}
\newcommand{\im}{\operatorname{im}}
\newcommand{\id}{\operatorname{id}}
\newcommand{\Sym}{\operatorname{Sym}}
\newcommand{\Gm}{\mathbf G_m}
\title{Hodge groups, exceptional classes, and odd-exponent abelian covers}
\author{Kyle Kistner}
\date{5 September 2026}
\begin{document}
\maketitle
\begin{abstract}
We prove a transfer criterion for full branch families of abelian covers
of $\mathbf P^1$ of any odd exponent $N$: the rational Hodge conjecture
on every power of the degree-$N$ Fermat Jacobian implies it on every
power of the very general covering Jacobian. Aoki's theorem makes this
unconditional for odd prime powers and for $N=3^a5^b$ or $3^a7^b$.
The proof includes the exceptional constant primitive factor at
conductor fifteen and its algebraic elliptic realization.
For prime-exponent families, actual oriented graph correspondences determine a rational
ambient group whose derived subgroup is connected monodromy; a signed
determinant lattice then gives the exact generic Hodge group. Branch
words formed from an odd multiplicative subgroup and opposite pairs yield
moving absolutely simple Jacobians with an explicitly counted first
exceptional Hodge quotient. These classes are algebraic already on the
Jacobian itself and generate all additional Hodge classes on its powers
under homomorphism pullbacks. The Hodge rings of distinct-prime parts
factor even when their covers share the full branch configuration. Finally,
the degree-seven order-three family satisfies the generalized Hodge conjecture
on every power, as an explicit corollary of classical cyclic monodromy,
Fermat-cycle algebraicity, and Abdulali's domination theorem. These are
written arguments depending on the stated companion results and standard
published inputs. They are not complete Lean formalizations, independently
refereed results, or claims of priority.
\end{abstract}

\section{Scope and dependencies}
The companion \cite{General} treats a full ordered branch family of connected
elementary covers of $\mathbf P^1$ of fixed prime exponent. It proves the
connected monodromy blocks, constructs actual oriented graph correspondences,
and realizes their invariant tensors through rational algebraic maps from
determinant-word Hodge structures. The companion
\cite[Corollary 1.2]{Fermat} transfers these determinants to prime-degree
Fermat varieties and applies Shioda's published product theorem. It proves
algebraicity of every rational Hodge class in their tensor products, including
all-row-balanced mixed determinants with repeated curves and different chosen
generators. The compact-type fusion argument \cite{Fusion} supplies an
alternative route for the balanced subspaces. We use these specific results,
not the frozen earlier manuscripts.

Section~\ref{sec:odd-exponent} supplies the additional primitive-factor
and constant-factor arguments needed at odd composite exponents. Its
cycle premise is the all-powers conjecture for one Fermat Jacobian,
verified in the stated cases by \cite{Aoki2000}; the Shioda--Katsura
construction \cite{SK} transfers that premise to the required Fermat
products. The earlier sections retain their stated prime-exponent scope.

Write $\Hdg^d(A)=H^{2d}(A,\Q)\cap H^{d,d}(A)$, and let $\Div^d(A)$
be the degree-$2d$ part of the rational algebra generated by divisors.
``Exceptional'' below means outside that algebra. All assertions concern
rational cohomology over $\C$. Very general always refers to the indicated
full configuration space, not an arbitrary restricted subfamily.

\section{The exact generic Hodge group}\label{sec:hg}
Fix an odd prime $p$, $K=\Q(\zeta_p)$, and
$\Gamma=\F_p^\times=\Gal(K/\Q)$. Let $G=(\Z/p)^m$ and let
$a_1,\ldots,a_r\in G\setminus\{0\}$ generate $G$ with sum zero.
At a very general point of the full ordered branch family put
$A=J(X)$ and $V=H^1(A,\Q)$. Use the evaluation code
\[
 C=\{(c(a_i))_i:c\in\Hom(G,\F_p)\}\subset\F_p^r.
\]
For a nonzero word $c$, put $s(c)=|\{i:c_i\ne0\}|$ and
$h(c)=s(c)-2$. Discard support-two factors, which have genus zero.
Over $K$ write $V_c$ for the character space; its dimension is $h(c)$.
All spaces retained have support at least three.

Define oriented classes as follows. On a support of size four, $c$ and $e$
are equivalent when $e=Pc$ for a labelled Klein-four double transposition
or the identity, with both words in $C$. On other supports the classes are
singletons. Write $O_c$ for the class. Negation acts on these classes.
By \cite[Section 3]{General}, each oriented relation comes from an actual
rational algebraic graph map on the cyclic factors, after one common finite
base cover. A self-negative class occurs only in rank two, has
opposite-paired residues, and has signature $(1,1)$ at every embedding.

Let $\mathcal E\subset\End_\Q(V)$ be the algebra generated by the deck
action, the oriented graph maps transported to $A$ by quotient pullback and
norm, and their polarization adjoints. Define
\[
 L=\Cent_{\Sp(V)}(\mathcal E)^0,
 \qquad M=\text{connected algebraic monodromy on }V.
\]
Every generator of $\mathcal E$ is an algebraic Hodge endomorphism, and
commutes with monodromy after the stated finite cover.

\begin{lemma}\label{lem:ambient}
The group $L$ is split over $K$, and
\[
 L_K\simeq
 \prod_{\{O,-O\},\ O\ne-O}\GL_{h(O),K}
 \ \times\prod_{O=-O}\SL_{2,K}.
\]
A non-self-negative pair contributes standard occurrences on $O$ and dual
occurrences on $-O$. Support-three classes contribute $\GL_1$ factors.
Moreover $L^{\mathrm{der}}=M$ as subgroups of $\Sp(V)$.
\end{lemma}
\begin{proof}
The character idempotents
\[
 e_c=\frac1{|G|}\sum_{g\in G}\zeta_p^{-c(g)}g^*
       \in K[G]
\]
are defined over $K$. Commuting with the deck action therefore forces
preservation of every $V_c$ already over $K$. The polarization pairs
$V_c$ nondegenerately with $V_{-c}$. If $\phi$ is an actual rational
oriented graph map, its component $e_e\phi e_c$ is a $K$-linear
isomorphism $V_c\to V_e$. Centralizing it requires conjugate actions on
these two spaces. Different paths between the same spaces differ by a
scalar: they commute with the irreducible standard monodromy action;
this can first be checked over $\C$ and descended to $K$. Thus additional
paths impose no further equations on a chosen matrix.

For $O\ne-O$, the free matrix on a representative space is an arbitrary
$\GL_h(K)$ matrix; all other matrices are obtained by the graph maps and
polarized duality. In a self-negative class, an actual map
$V_c\to V_{-c}$ combined with the polarization gives a nondegenerate
monodromy-invariant bilinear form on $V_c$. The standard $\SL_2$
representation has its alternating form as its unique invariant bilinear
form, up to a scalar. The resulting form is therefore alternating over
$K$, and its preserving group is the split $\Sp_2=\SL_2$.

This calculates the centralizer over $K$, not merely its complexification;
all displayed factors are connected. Its derived subgroup is precisely
the product of the moving $\SL_h$ blocks in
\cite[Theorem 3.1]{General}. Support-three factors have trivial connected
monodromy. Hence its derived subgroup and $M$ are equal inside the original
faithful representation.
\end{proof}

Put $D=L/M$ and let $q:L\to D$ be the quotient. Choose one class from
each non-self-negative pair. Its determinant character is denoted $e_O$;
then
\[
 E=X^*(D_K)=\bigoplus_{\{O,-O\},\ O\ne-O}\Z e_O.
\]
Reversing orientation changes $e_O$ to $-e_O$. A self-negative class
contributes no generator; retaining a formal relation $2e_O=0$ would give
an incorrect unsaturated presentation. The determinant on each $\GL_h$
factor has kernel the actual $\SL_h$, so no $h$th power belongs in this
quotient map. This is the quotient-torus lattice, rather than an
identification with the lattice of an isogenous central torus.

Galois acts on $E$ by the signed permutation of oriented labels. The
$K$-rational calculation in Lemma~\ref{lem:ambient} justifies that these
are the complete descent data on $D$: changes of graph paths or bases act
by inner conjugation and leave determinants unchanged.

For a representative $c\in O$, set
\[
 d_c(b)=s(c)-1-\frac1p\sum_{i:c_i\ne0}[bc_i]_p,
 \qquad \delta_c(b)=2d_c(b)-h(c),
\]
where $[x]_p\in\{1,\ldots,p-1\}$. Oriented graphs preserve $d$,
negation changes the sign of $\delta$, and $\delta=0$ on self-negative
classes. Define the integral map
\[
 \Sigma:E\longrightarrow\Z^\Gamma,
 \qquad e_O\longmapsto(\delta_c(b))_{b\in\Gamma}.
\]
Its kernel is saturated. Let $T\subset D$ be the subtorus with annihilator
$\ker\Sigma$, so that
\[
 X^*(T_K)=E/\ker\Sigma\simeq\im\Sigma.
\]
The quotient is primitive; this does not claim that $\im\Sigma$ is
saturated in the unrelated ambient lattice $\Z^\Gamma$.

\begin{theorem}\label{thm:hg}
At a very general point of this full branch family,
\[
                  \Hg(A)=q^{-1}(T)\subset L.
\]
Consequently
\[
 \dim\Hg(A)=\sum_{\text{moving complex blocks}}(h^2-1)
                       +\operatorname{rank}_{\Q}\Sigma,
 \qquad \MT(A)=\Gm\cdot\Hg(A)\subset\GL(V).
\]
The sum counts each complex simple block once, including self-negative
rank-two blocks.
\end{theorem}
\begin{proof}
The generic monodromy inclusion gives $M\subset\Hg(A)$, while the
algebraic Hodge endomorphisms defining $L$ give $\Hg(A)\subset L$.
The projected Hodge circle in $D$ pairs with $e_O$ by the integer
$2d_c(1)-h(c)$. The smallest rational subtorus containing this circle is
obtained by adjoining all its Galois-conjugate cocharacters. Since $D$
splits over $K$, their pairings are exactly the rows of $\Sigma$.
Its annihilator is therefore $\ker\Sigma$, and this subtorus is $T$.

By minimality of the Hodge group, $\Hg(A)\subset q^{-1}(T)$; its image
under $q$ contains $T$ and is contained in $T$. Since it contains the full
kernel $M$, it equals $q^{-1}(T)$. This preimage is connected because
both its kernel and quotient are connected. The dimension and weight-one
Mumford--Tate formulas follow. The last product is not asserted to be a
direct product.
\end{proof}

\begin{remark}
The minimal rational span of a Hodge-signature vector is standard: compare
Xue--Zarhin \cite[Lemma 3.9, Theorem 3.11 and Section 3.13]{XZ}.
The geometric work here is the exact ambient centralizer with its actual
oriented maps, including self-negative rank-two and constant rank-one
sectors. An unproved replacement of Hodge endomorphisms by the entire
monodromy commutant would not establish this theorem.
\end{remark}

\section{Exceptional classes on a single simple Jacobian}\label{sec:exceptional}
Let $P\subset\F_p^\times$ be a multiplicative subgroup of odd order
$m\ge3$. Choose $k\ge0$ with $s=m+2k\ge5$, and take the ordered
branch word
\begin{equation}\label{eq:subgroupword}
 a=(\text{one occurrence of every element of }P,
                         \underbrace{1,-1,\ldots,1,-1}_{k\text{ pairs}}).
\end{equation}
Its sum is zero, because the elements of $P$ are the roots of $x^m-1$.
Let $Y\to\mathbf P^1$ be the associated cyclic cover and $A=J(Y)$.
Put
\[
 h=s-2=m+2k-2\ge3,\qquad g=\frac{(p-1)h}{2},\qquad r_0=(p-1)/2.
\]
The integer $h$ is odd. Write $V_b\subset H^1(A,\C)$ for the deck
eigenspace and $D_b=\bigwedge^hV_b$ for its volume line.

\begin{theorem}\label{thm:exceptional}
For every smooth member of this family there is an algebraic rational
subspace $\mathscr E\subset\Hdg^h(A)$ with
\begin{equation}\label{eq:exceptionalspace}
 \mathscr E_\C=
 \bigoplus_{\substack{\{b,-vb\}\text{ distinct unordered pairs}\\
             b\in\F_p^\times,\ v\in P\setminus\{1\}}}
             D_b\wedge D_{-vb},
 \qquad \dim_\Q\mathscr E=\frac{(p-1)(m-1)}2.
\end{equation}
At a very general member, $\End^0(A)=K$, the Jacobian is absolutely simple
and not of CM type, and
\[
 \Hdg^d(A)=\Div^d(A)\quad(d<h),\qquad
 \Hdg^h(A)=\Div^h(A)\oplus\mathscr E.
\]
In particular $h$ is its first exceptional codimension, and
\[
 \dim_\Q\Hdg^h(A)=
        \binom{r_0+h-1}{h}+\frac{(p-1)(m-1)}2.
\]
\end{theorem}

\begin{proof}[Algebraicity and diagonal realization]
Fix $v\in P\setminus\{1\}$. Use $Y$ twice, with second generator
$\sigma_2=\sigma^{-v^{-1}}$. A generator change $\sigma\mapsto\sigma^e$
changes residues by $a_i\mapsto e^{-1}a_i$. Thus the second residue word
is $-va$, and its $b$-eigenspace is the original $V_{-vb}$.
The concatenation $a\sqcup(-va)$ is a union of opposite pairs:
its subgroup portion is $P\sqcup(-P)$ because $vP=P$, and its added
pairs remain opposite pairs. The cyclic Hodge-dimension formula therefore
makes the following rational space entirely of Hodge type $(h,h)$:
\[
 (W_v)_\C=\bigoplus_{b\ne0}D_b\otimes D_{-vb}
                    \subset H^h(A)\otimes H^h(A)
\]
It has dimension $p-1$ and is algebraic by the Fermat transfer
\cite[Corollary 1.2]{Fermat}. Equivalently one can use balanced fusion
\cite[Theorem 1.1]{Fusion}.

Pull back along $\Delta:A\to A\times A$. Since $P$ has odd order,
$-1\notin P$ and $b\ne-vb$; cup multiplication is nonzero on every
displayed line. Equality $\{b,-vb\}=\{b',-vb'\}$ either gives
$b=b'$ or, by exchanging endpoints, $v^2=1$. The latter is impossible
for $v\ne1$ in $P$. Hence $\Delta^*$ is injective on $W_v$, and its
image $\mathscr E_v$ has dimension $p-1$.

For different parameters, endpoint exchange shows that their unordered
label pairs coincide exactly when the parameters are $v$ and $v^{-1}$.
Thus $\mathscr E_v=\mathscr E_{v^{-1}}$, and otherwise their exterior
multi-degrees are disjoint. Summing over the $(m-1)/2$ inverse pairs
proves \eqref{eq:exceptionalspace}. Rationality and algebraicity follow
from the rational determinant spaces and algebraic diagonal pullback.
\end{proof}

\begin{proof}[Simplicity and the complete first quotient]
For this cyclic family $s\ge5$, the connected monodromy theorem gives
\[
 M_\C=\prod_{\{b,-b\}}\SL_h,
\]
standard on $V_b$ and dual on $V_{-b}$. Since $h\ge3$, these two
representations are inequivalent. Its commutant on $H^1(A,\C)$ consists
of a scalar on each of the $p-1$ eigenspaces. The known rational algebra
$K$ has this full dimension, so $\End_M(H^1(A,\Q))=K$. At a very
general point every Hodge endomorphism commutes with $M$; the algebraic
deck action then gives $\End^0(A)=K$. A proper abelian subvariety would
give a nontrivial rational idempotent by Poincare reducibility, impossible
in this field. Also $[K:\Q]=p-1<2g$, so $A$ is not CM.

For one block $U\oplus U^*$ with $\dim U=h$, the exterior invariants
of $\SL_h$ are generated by a polarization tensor $\theta$ of degree
two and the two volume lines of degree $h$. Indeed
$\bigwedge^iU$ is irreducible for $0<i<h$, and
\[
 (\bigwedge^iU\otimes\bigwedge^jU^*)^{\SL_h}
\]
is one-dimensional for $i=j$ and otherwise zero, apart from the trivial
endpoint exterior powers for $i,j\in\{0,h\}$. The paired invariants
are the nonzero powers $\theta^i$, $0\le i\le h$, and the product of
the two volumes is a nonzero multiple of $\theta^h$.

Invariants for the product group tensor together. In even degree below
$2h$, the oddness of $h$ prevents a single volume, while two volumes have
too large a degree. Only products of pairings remain. The degree-two
invariant space consists entirely of divisors: equivalently
\[
 \operatorname{NS}(A)_\Q\simeq\End^0(A)^{\dagger=1}=K^+,
\]
and its complexification contains each of the $r_0$ pairing lines.
Thus all Hodge classes below codimension $h$ are divisor-generated.

In degree $2h$, the additional invariants are pairs of volumes.
Opposite labels already give powers of a pairing, and equal labels give
zero in the exterior algebra. The remaining rational Galois orbits are
$\mathscr E_v$ for $v\ne\pm1$, with $v$ identified with $v^{-1}$.
The diagonal argument above gives each such orbit a rank-one $K$-Hodge
structure of dimension $p-1$, regardless of whether it is Hodge.
Its rational Hodge subspace is zero or the whole space. Writing
$d_a(b)=\dim V_b^{1,0}$, it is the whole space exactly when
\begin{equation}\label{eq:stabilizer}
                   d_a(vb)=d_a(b)\quad\text{for every }b\ne0.
\end{equation}

The stabilizer in \eqref{eq:stabilizer} is exactly $P$. Inclusion of $P$
follows from its multiplicative invariance and the constant contribution
of every opposite pair to the fractional-part sum. Conversely
\eqref{eq:stabilizer} is all-row balance for $a\sqcup(-va)$.
The prime balanced-tuple equality of Aoki--Parry, in the precise form
recorded in \cite[Section 3]{Fusion}, forces its residue multiplicities
to be invariant under negation. Cancelling the added pairs gives
\[
                 1_P-1_{-P}=1_{vP}-1_{-vP}.
\]
Distinct multiplicative cosets are disjoint and $P\ne-P$, so this
equality forces $vP=P$. Thus precisely the spaces in
\eqref{eq:exceptionalspace} supply the additional Hodge classes.

Divisor products have equal occupation numbers in $V_b$ and $V_{-b}$
for every label pair. A line $D_b\wedge D_{-vb}$ with $v\ne\pm1$
does not. Their exterior multi-degrees are disjoint, proving
$\mathscr E\cap\Div^h(A)=0$. The pairing algebra is
\[
 \C[\theta_1,\ldots,\theta_{r_0}]/(\theta_i^{h+1}).
\]
In degree $h$ in these generators the truncations remove no monomial,
so its dimension is $\binom{r_0+h-1}{h}$. This proves the assertion.
\end{proof}

\begin{remark}[Moving and unbounded examples]
The parameter space has dimension $s-3>0$. A fixed curve of genus at
least two has finitely many cyclic actions of order $p$, so forgetting
the covering structure has finite fibers; Torelli gives the same dimension
for the Jacobian image.

For every prime $p\equiv1\pmod3$, choose an order-three element $u$ and
take $P=\{1,u,u^2\}$, $k\ge1$. Then $h=2k+1$,
$g=(p-1)(2k+1)/2$, and $\dim\mathscr E=p-1$.
At $p=7$, the word $(1,2,4,1,6)$ gives a two-dimensional family of
absolutely simple ninefolds with
\[
                 \dim\Hdg^3(A)=10+6=16.
\]
The six exceptional lines have unordered labels
$\{1,5\},\{2,3\},\{1,3\},\{4,6\},\{4,5\},\{2,6\}$.

Alternatively, for $p\equiv3\pmod4$, $p\ge11$, use the quadratic
residues and $k=0$. Then
\[
 h=(p-5)/2,\quad g=(p-1)(p-5)/4,\quad
 \dim\mathscr E=(p-1)(p-3)/4.
\]
The word $(1,3,4,5,9)$ at $p=11$ gives twenty exceptional algebraic
codimension-three classes on a moving simple fifteenfold. These formulas
give unbounded prime degrees, dimensions and codimensions.
\end{remark}

\section{Generators on every power of the subgroup family}\label{sec:ring}
The first exceptional space suffices on all powers, provided it is allowed
to move between copies by homomorphism pullback.

\begin{theorem}\label{thm:ring}
For a very general $A$ in \eqref{eq:subgroupword}, and every $q\ge1$,
the rational Hodge ring of $A^q$ is generated by its divisors and
\[
       f^*\mathscr E\subset\Hdg^h(A^q),\qquad
                       f\in\Hom(A^q,A).
\]
All these generators are algebraic. No exceptional Hodge class occurs
on any power in codimension strictly less than $h$.
\end{theorem}
\begin{proof}
The rational diagram and Hodge-lifting proof of
\cite[Sections 5--6]{General} expresses every Hodge class as a span of
pairing diagrams and all-row-balanced effective volume words. For this
cyclic family an effective word is a list of labels $b_1,\ldots,b_N$;
its branch words are the scalar multiples $b_j a$. The added opposite
pairs make no contribution to antisymmetric residue multiplicities.
Aoki--Parry therefore says that balance is equivalent to
\[
     \sum_j (1_{b_jP}-1_{-b_jP})=0.
\]
As $P\ne-P$, this means that each coset and its negative occur equally
often. Pair the volume occurrences accordingly. Each pair has labels
$b,-vb$ for some $v\in P$. This pairing of a word is compatible with
its full Galois orbit: simultaneous multiplication of all labels preserves
each chosen value of $v$. No negative multiplicities are introduced.

For $v=1$, the paired volume tensors are alternating products of
standard--dual coevaluations; hence their realizations are generated by
divisors. For $v\ne1$ they are the volume pairs defining $\mathscr E$.
It remains to show that homomorphism pullbacks realize all their possible
placements in the copies of $A^q$, rather than just placing both volumes
in one copy.

For a label $b$, the standard space in $H^1(A^q,\C)$ is
$V_b\otimes\C^q$. Its alternating $h$-volume subspace is
\[
                  D_b\otimes\Sym^h(\C^q).
\]
This follows either from the exterior Cauchy formula or by alternating
the $V_b$ slots: because their determinant is alternating, the copy
indices are symmetric. Pullback by
$f=\sum_i a_i\operatorname{pr}_i$, $a_i\in K$, sends $D_b$ to its
tensor with the pure $h$th power of the coefficient vector
$(\tau_b(a_i))_i$. For a nonopposite pair $b,c=-vb$, the vectors at
$b$ and $c$ vary independently after taking complex linear spans.
Indeed $K^q$ is Zariski dense in the affine $\Q$-space
$\operatorname{Res}_{K/\Q}\mathbb A^q$: choose a $\Q$-basis of $K$
and use the density of $\Q^{q(p-1)}$. After extension to $\C$ the
embedding coefficient vectors are independent coordinates. Pure powers
span symmetric powers in characteristic zero. Thus these pullbacks span
\[
 (D_b\wedge D_c)\otimes
                  \Sym^h(\C^q)\otimes\Sym^h(\C^q),
\]
including every placement of this pair of volumes.

Rational coefficients in the $a_i$ cause no problem: clearing a common
denominator gives an actual homomorphism, and its action in degree $2h$
differs by a nonzero rational scalar. Likewise use a rational basis of
$\mathscr E$ before extending scalars; no single complex line is assumed
to be a rational generator. Cup products of the resulting pair-volume
classes and divisors span each balanced diagram in its prescribed copy
positions. Faithful scalar extension then gives the asserted rational
generation. Every generator is algebraic by Theorem~\ref{thm:exceptional}
and algebraic pullback. Finally a tensor invariant of even degree below
$2h$ cannot contain an odd number of degree-$h$ volumes, nor two such
volumes; only the divisor pairings remain.
\end{proof}

\section{Products belonging to different primes}\label{sec:products}
The distinction between separate validity of the Hodge conjecture and its
validity for a product requires an additional argument. Here independent
cyclotomic embeddings provide it, including for constant CM factors.

\begin{lemma}\label{lem:fields}
Let $K_1,\ldots,K_a$ be jointly linearly disjoint finite Galois CM fields.
Let $W_i$ be a pure rational Hodge structure of weight $n_i$, rank one
over $K_i$, with its field action by Hodge endomorphisms. Suppose
$\sum n_i+2t=2d$. Then
\[
                  \bigotimes_i W_i\otimes\Q(-t)
\]
has a nonzero rational Hodge class if and only if every $n_i$ is even
and $W_i$ is entirely of type $(n_i/2,n_i/2)$. In that event its Hodge
vectors are spanned by external tensors of individual Hodge vectors and
the Tate vector.
\end{lemma}
\begin{proof}
The tensor product $K_1\otimes\cdots\otimes K_a$ is the compositum
field, and the displayed source has rank one over it. Its rational Hodge
subspace is field-stable, so is zero or the whole source. The embeddings
of this field are independent tuples of embeddings of the $K_i$.
If $f_i(\tau)$ is the first Hodge index on the $\tau$-line of $W_i$,
all tuple sums $\sum f_i(\tau_i)+t$ must therefore equal $d$.
Varying one embedding makes each $f_i$ constant. Conjugation gives
$f_i(\bar\tau)=n_i-f_i(\tau)$, forcing the constant to be $n_i/2$.
The converse and the tensor-span assertion follow immediately.
\end{proof}

For distinct odd primes, the fields $K_p=\Q(\zeta_p)$ are jointly
linearly disjoint: their compositum is $\Q(\zeta_{\prod p})$ and its
degree is $\prod(p-1)$, with Galois group $\prod\F_p^\times$, by
the cyclotomic degree theorem and the Chinese remainder theorem.
The binary part has only Tate invariant sources.

\begin{theorem}\label{thm:products}
For each prime in a finite set, prescribe a connected elementary cover of
$\mathbf P^1$ with that prime exponent. Either let their full ordered
branch configurations vary independently, or use one common full ordered
configuration, allowing arbitrary overlaps among active labels. At a
very general parameter, write $A_p$ for their Jacobians. For every choice
of nonnegative powers $q_p$, external product gives an isomorphism
\[
 \bigotimes_p\Hdg^*(A_p^{q_p})
                  \ \simeq\ \Hdg^*\left(\prod_p A_p^{q_p}\right).
\]
Every class on the right is algebraic. Each $A_p$ includes its constant
CM part. The assertion also passes to abelian isogeny factors of these
prime parts.
\end{theorem}

\begin{proof}[Independence of the moving groups]
For independently varying configurations, the fundamental group of the
product parameter space is the product and the monodromy acts in separate
summands. Its connected closure is the product of the individual groups.
Consider a common full configuration instead. Forgetting inactive labels
has connected configuration-space fibers, giving full projection on each
individual monodromy group. The connected group is semisimple: its radical
projects trivially to every semisimple moving projection and acts
trivially on the constant factors. A proper subdirect product would, by
semisimple Goursat, identify two isomorphic simple adjoint blocks.

Fix an odd-prime moving character of prime $p$. There is an active pair
$i,j$ whose residue sum is nonzero; otherwise three active nonzero entries
would have all pair sums zero, contradicting that $p$ is odd. The labelled
pair twist has eigenvalues
$(\zeta_p^{c_i+c_j},1,\ldots,1)$ projectively, and hence exact
projective order $p$. On a different odd-prime $\ell$ character, the same
twist is either identity (an inactive label), unipotent (zero pair sum),
or has projective order $\ell$. These are the pair formulas and
nonzero-rank-one checks in \cite[Section 3]{General}, from
Menet--Nguyen \cite[Theorem 2.5]{MN}.

This comparison is legitimate before finite auxiliary base changes. The
projective representations are defined on the original common marked-sphere
fundamental group. A graph between two centerless simple groups is its
own normalizer in their product: a normalizing pair differs from a graph
pair by an element centralizing the target group. Since connected closure
is normal in full closure, any connected graph relation would hold on
every original pair twist. An algebraic group automorphism preserves
element order, so the orders just listed exclude such a graph across
distinct odd primes.

An odd-prime $\SL_h$ and binary $\Sp_{2g}$ can have isomorphic simple
groups only for $h=2,g=1$: equal rank and dimension give
$h-1=g$ and $h^2-1=g(2g+1)$, hence $g(g-1)=0$.
In that case the same binary twist is identity or a
Picard--Lefschetz unipotent, of infinite order when nonidentity; it cannot
match the odd-prime finite order. Thus all cross-prime graphs are excluded.
The connected monodromy is exactly the product of the prime-part groups.
\end{proof}

\begin{proof}[Hodge-source separation and algebraicity]
In every cohomological degree $n$, the rational diagram construction of
\cite[Section 5]{General} gives an algebraic rational Hodge surjection
onto $H^n(A_p^{q_p},\Q)^{M_p}$ from a finite sum of sources
\[
                        W_{p,w}\otimes\Q(-t),
                        \qquad H(w)+2t=n,
\]
and Tate sources. Each nonempty word source is rank one over $K_p$.
Although the companion's final application sets $n=2d$, the diagram
construction itself works in arbitrary degree: none of its slot,
alternation, Galois-orbit or graph operations requires even $n$.
This observation is necessary for the odd/odd K\"unneth components here.
For $p=2$ the invariant sources are entirely Tate.

In each K\"unneth multidegree, independence of the moving groups makes
the invariant target the tensor product of its individual invariant
spaces. Tensor their source maps. The resulting sources are finite sums
of tensor products with at most one rank-one $K_p$ word for each odd
prime, together with Tate twists. Lemma~\ref{lem:fields} says that a
nonzero Hodge vector in such a source forces every individual prime-word
source to be balanced and of even weight. In particular odd prime-part
cohomological degrees supply no Hodge classes in the global decomposition.

At a very general point, Hodge classes are fixed by connected monodromy.
Lift each through this rational Hodge surjection using semisimplicity of
polarizable pure Hodge structures. Its components are Hodge separately
in the displayed rational source summands, so they are spanned by external
products of individual balanced Hodge vectors. The forward maps are
algebraic and each balanced word is algebraic by
\cite[Corollary 1.2]{Fermat}, or alternatively by same-prime fusion.
This proves algebraicity and external-product surjectivity. K\"unneth
gives injectivity and compatibility with multiplication.

No algebraicity of the Hodge splitting is needed. Support-three constant
factors are retained as rank-one determinant entries throughout.
An abelian isogeny factor is reached by rational algebraic inclusion and
retraction; pulling back and then retracting classes proves the same
assertions for its powers. Prime-part algebraic projectors also preserve
the stated factorization.
\end{proof}

\begin{remark}[Product boundary]
Theorem~\ref{thm:products} concerns products of prime-cover Jacobians
and their isogeny factors; the entire Jacobian of a composite-exponent
compositum has additional mixed-order characters. The separate criterion
in Section~\ref{sec:odd-exponent} treats those full Jacobians at its
stated odd exponents. Neither assertion replaces the full configuration
by a subfamily. Jointly very general
parameters are required; separate genericity alone is not substituted.
The special Hodge groups can equivalently be recovered as a product from
the factorization of all tensor invariants, but the full Mumford--Tate
groups share their weight direction and must not be described as an
unrestricted direct product.
\end{remark}

\section{A classical generalized-Hodge corollary in degree seven}
\label{sec:ghc7}
Here the stronger conjecture follows from an additional domination argument,
not from ordinary all-powers algebraicity alone. We give the argument because
the auxiliary CM varieties can be identified explicitly. The ordinary-Hodge
input in this section uses only a single moving cyclic cover and a fixed
CM elliptic curve; the simultaneous elementary-abelian theorem is unnecessary.

\begin{theorem}\label{thm:ghc7}
For $k\geq1$, let $A$ be the Jacobian at a very general point of the full
ordered cyclic degree-seven branch family with word
\[
                    (1,2,4,(1,-1)^k).
\]
Then $\dim A=3(2k+1)$, and the generalized Hodge conjecture holds for
every power of $A$.
\end{theorem}

\begin{proof}[Signatures and the domination hypothesis]
Put $h=2k+1$, $K=\Q(\zeta_7)$, $V=H^1(A,\Q)$, and
$P=\{1,2,4\}\subset\F_7^\times$.
Rohde's prime-cyclic monodromy theorem
\cite[Theorem 5.3.7]{Rohde} gives $M_\C=\SL_h^3$.
As in Section~\ref{sec:exceptional}, its commutant is $K$, so
$\End^0(A)=K$ and $A$ is absolutely simple of Albert type IV.
The inclusions $M\subset\Hg(A)\subset L(A)$, where $L(A)$ is the
Lefschetz group, give
\[
                   \Hg(A)^{\mathrm{der}}=L(A)^{\mathrm{der}}=M.
\]
This is precisely the definition of pel-type in \cite[Section 2]{Abdulali12};
it does not assert that the branch family fills a PEL moduli space.

Multiplication by $b\in P$ permutes $(1,2,4)$, whereas its negative has
least-residue sum fourteen. Each added opposite pair has sum seven.
Thus the Hodge multiplicities of $V_b$ are $(k+1,k)$ for $b\in P$
and $(k,k+1)$ for $b\in-P$. Every real unitary signature differs by
exactly one. Abdulali's Theorem 10 therefore makes $A$ weakly
self-dominated, and his Lemma 9 gives the same property for every
$A^q$. The difference-one condition is essential: it is also the
correction to the older type-IV statements in
\cite[Appendix B.5]{AbdulaliSurvey}. We use the 2012 formulation with
that condition explicitly imposed.
\end{proof}

\begin{proof}[Identification of every required CM partner]
Let $D=\bigwedge_K^h V$ denote the rational determinant realization.
Its twist $D(k)$ has weight one, with type $(1,0)$ on $P$ and $(0,1)$
on $-P$. The three-point curve
\[
                          C_0:\quad y^7=x(x-1)^2
\]
has these same $K$-types. Both spaces have rank one over $K$, so
$D(k)\simeq H^1(C_0,\Q)$ as $K$-linear Hodge structures.

The fixed field $K^P$ is $K_0=\Q(\sqrt{-7})$. Indeed,
$z=\zeta_7+\zeta_7^2+\zeta_7^4$ satisfies
$z+\bar z=-1$ and $z\bar z=2$, hence $z^2+z+2=0$.
Choose a CM elliptic curve $E$ with field $K_0$ and the type whose
extensions to $K$ are $P$. Its induced $K$-Hodge structure has the
preceding types, whence
\[
 D\simeq H^1(E,\Q)^{\oplus3}(-k),
 \qquad J(C_0)\sim E^3.
\]
Here the isogeny follows from the equivalence of rational weight-one
Hodge structures and complex abelian varieties up to isogeny.

It remains to handle the entire determinant tensor algebra, not just
$D$. The Hodge group of $E$ is the quadratic norm-one torus. For
$n\geq1$ its rational character pair $\chi^n\oplus\chi^{-n}$ has
a realization $U_n$ of types $(n,0),(0,n)$ in $H^n(E^n,\Q)$: take
the rational sum of the all-holomorphic and all-antiholomorphic lines
in $H^1(E)^{\otimes n}$. The two lines are Galois conjugate, so their
sum is rational. Put $U_0=\Q\subset H^0(E,\Q)$.
Every irreducible constituent of $D^{\otimes N}$ is a copy of
\[
 U_n\left(-Nk-\frac{N-n}{2}\right),
 \qquad 0\leq n\leq N,\quad n\equiv N\pmod2.
\]
Twisting by $Nk+(N-n)/2$ realizes it as $U_n$ in a power of $E$.
These realizations are fully twisted: they are effective and have a
nonzero form of type $(n,0)$. Trivial constituents use $U_0$; duals
and Tate factors change only the twist, not the CM character.

In the proof of \cite[Theorem 11]{Abdulali12}, the auxiliary CM
category is the tensor algebra of the determinants of the simple
isogeny factors. The only simple factor of $A^q$ is $A$, so the
object just computed is exactly the required one. By
\cite[Remark 12]{Abdulali12}, every $A^q$ is consequently dominated
by varieties $A^n\times E^t$: each irreducible Hodge structure in
its cohomology has a fully twisted Tate-equivalent realization on
one of these varieties. No assertion about arbitrary CM types over
$K$ is involved.
\end{proof}

\begin{proof}[Ordinary Hodge conjecture on the required products]
We verify ordinary HC for all $A^u\times E^v$ directly from the
classical cyclic and Fermat inputs. The fixed factor $E$ has trivial
monodromy over the branch configuration, so connected monodromy of
the product remains $M$, acting only on $V$.
Special-linear invariant theory expresses all $M$-invariant tensors
through polarization contractions, full determinant tensors, and
unrestricted tensors in $H^1(E)$. Taking rational Galois orbits of
the diagrams gives a rational Hodge surjection, in each degree,
from a finite direct sum of tensor products of copies of $D$,
$H^1(E)$, and Tate factors onto the monodromy invariants.
Its maps are algebraic correspondences: deck projectors, permutations,
alternation, polarization, and the usual cup-product correspondences.
This is the single-cyclic instance of the explicit construction in
\cite[Section 5]{General}; it requires no comparison of distinct
moving cyclic quotient lines.

The determinant $D$ is an algebraic direct summand of the middle
cohomology of a degree-seven Fermat $h$-fold, by the classical
finite-quotient construction written out in
\cite[Theorem 1.1]{Fermat}. The curve $C_0$ is a quotient of the
degree-seven Fermat curve: in affine coordinates $r^7+s^7=1$, set
$x=r^7$ and $y=rs^2$. Then $y^7=x(x-1)^2$.
Since $E$ is an isogeny factor of $J(C_0)$, its first cohomology is
an algebraic direct summand of that of the Fermat curve. Thus every
source just described is an algebraic direct summand of the
cohomology of a product of degree-seven Fermat varieties, with
Tate factors supplied by projective spaces.

Shioda's fixed-prime product theorem \cite[Theorem 2]{Shioda79}
gives ordinary HC on all these Fermat products; projective-space
factors preserve it by the projective bundle formula. At a very
general point, each Hodge class of $A^u\times E^v$ is fixed by $M$.
Semisimplicity of polarizable rational Hodge structures supplies a
Hodge preimage under the displayed surjection. The algebraic source
projectors and Shioda make that preimage algebraic, and the forward
diagram map transports its class to $A^u\times E^v$.
Only the forward map, not a chosen Hodge splitting, needs to be
algebraic. Taking the countable intersection over $u,v$ leaves one
very-general locus on which all these ordinary-Hodge assertions hold.
\end{proof}

\begin{proof}[Conclusion of Theorem~\ref{thm:ghc7}]
Fix $q\geq1$. The dominating varieties found above are
$A^n\times E^t$, and ordinary HC holds for
$A^q\times(A^n\times E^t)=A^{q+n}\times E^t$.
The domination criterion \cite[Proposition 1]{Abdulali12}, attributed
there to Grothendieck, proves generalized HC for $A^q$.
This holds for every $q$.
\end{proof}

\begin{remark}[Scope and attribution of the corollary]
This is an explicit consequence of Rohde's cyclic monodromy theorem,
the classical Fermat determinant construction and Shioda's cycle
theorem, and Abdulali's domination machinery. No novelty is claimed
for the resulting family. The explicit quadratic CM calculation
explains why the stronger conjecture follows here: it supplies all
the auxiliary fully twisted realizations required by the domination
proof. For other primes, realizing one determinant as a Fermat
Hodge structure does not by itself identify fully twisted
realizations for every irreducible constituent of its tensor
algebra. We make no analogous general-prime assertion.
\end{remark}

\section{All powers for odd-exponent abelian covers}
\label{sec:odd-exponent}
This section extends the prime-exponent results to full Jacobians of
abelian covers with odd composite exponent. The cycle-theoretic premise
is isolated in a fixed Fermat Jacobian, including all its powers.

\begin{theorem}[Odd-exponent transfer criterion]
\label{thm:odd-exponent-criterion}
Let $N>1$ be odd and assume the rational Hodge conjecture for every
power of the Jacobian of the degree-$N$ Fermat curve. Let $G$ be a
finite abelian group of exponent dividing $N$, and let
$a_1,\ldots,a_r\in G\setminus\{0\}$ generate $G$ and have sum zero.
Then the rational Hodge conjecture holds for every power of the
very general Jacobian in the full ordered branch family of connected
$G$-covers of $\mathbf P^1$ with these local monodromies.
\end{theorem}

Very general means outside a countable union of proper closed algebraic
subsets of $M_{0,r}$. A finite base cover may be used to choose the family
and auxiliary algebraic graphs. The case $r=2$ is a cyclic genus-zero
cover; $r=3$ is included as a constant family. The following subsections
prove the theorem and then identify unconditional exponents.

\subsection{Primitive quotient factors, without double counting}
Fix $\zeta_N$ and identify $\widehat G$ with
$\Hom(G,\Z/N)$. Evaluation on the branch elements gives an injective
code $C\subset(\Z/N)^r$; freeness of this module is not assumed.
Here $N$ is odd, and the exponent of $G$ divides $N$.
For $c\in C\setminus\{0\}$, let $d=d(c)$ be its additive order and
write
\[
 c_i=(N/d)b_i\quad\text{in }\Z/N,
 \qquad b_i\in\Z/d.
\]
The residues $b_i$ generate $\Z/d$, have sum zero, and have the same
active support $S(c)$ as $c$. Set $s(c)=|S(c)|$ and $h(c)=s(c)-2$.
The cyclic quotient $C_c=X/\ker c$ has degree $d$. Choose its deck
generator $T_c$ so that the specified character has eigenvalue
$\zeta_d$. Let
\[
 U_{[c]}=\ker\Phi_d(T_c^*:H^1(C_c,\Q)\longrightarrow H^1(C_c,\Q)).
\]
Here $[c]$ is the orbit under $(\Z/N)^\times$ and $\Phi_d$ is the
$d$th cyclotomic polynomial. The restriction map
$(\Z/N)^\times\to(\Z/d)^\times$ is surjective. Thus
\[
 U_{[c]}\otimes\C=\bigoplus_{u\in(\Z/d)^\times}V_{uc},
 \qquad \dim_{\Q(\zeta_d)}U_{[c]}=h(c).
\]
Taking one representative per nonzero orbit gives the rational
decomposition of $H^1(X,\Q)$, omitting the zero-dimensional support-two
terms. In particular, we do \emph{not} sum the entire cohomology of all
cyclic quotients: that would count their lower-conductor parts repeatedly.

All the displayed primitive projections are rational algebraic. For
any $d$, the projector on primitive eigenvalues of an order-$d$ operator is
\begin{equation}\label{eq:ppprimitive}
 e_d(T)=\sum_{k\mid d}\mu(d/k)\frac{k}{d}
                  \sum_{j=0}^{d/k-1}T^{kj},
\end{equation}
where $\mu$ is the M\"obius function. The inner average projects onto
eigencharacters whose orders divide $k$, so M\"obius inversion selects
order exactly $d$. For $d=p^f$ this simplifies to
$1-p^{-1}\sum_{j=0}^{p-1}T^{j d/p}$.
Composing this polynomial with quotient pullback and norm
realizes the primitive factor inside $H^1(X,\Q)$. Its polarization and
its cyclotomic action are consequently algebraic as well.

\subsection{Primitive monodromy at every odd conductor}
Let $d>1$ be odd and let $a=(a_1,\ldots,a_s)$ have nonzero entries in
$\mathbb Z/d$, with $\sum_i a_i=0$ and
$\gcd(d,a_1,\ldots,a_s)=1$. We consider the full ordered branch family
of cyclic degree-$d$ covers of $\mathbb P^1$ with this word. Put $h=s-2$.
For a unit $b\bmod d$, set
\[
 Q_a(b)=\frac1d\sum_i[ba_i]_d,\qquad
 (h^{1,0}_b,h^{0,1}_b)=(s-1-Q_a(b),Q_a(b)-1),
\]
where $[x]_d\in\{1,\ldots,d-1\}$ for $x\ne0$.
These are the multiplicities of the two weight-one Hodge types,
not the bidegrees of a weight-$h$ determinant.

\begin{lemma}[A good triple outside the small conductors]
\label{lem:odd-good-triple}
Suppose $s\ge4$ and $d\notin\{3,5,15\}$. At every primitive embedding,
either the eigenspace is mixed or there are three active labels $i,j,k$
such that the order of $\zeta_d^{a_i+a_j}$ exceeds $5$ and at least one
of $\zeta_d^{a_i+a_k},\zeta_d^{a_j+a_k}$ has order exceeding $2$.
Consequently each primitive complex eigenspace has connected complex
monodromy $\mathrm{SL}_h$.
\end{lemma}
\begin{proof}
If every pair eigenvalue had order at most $5$, its odd order would
divide $15$. Hence $d\mid15(a_i+a_j)$ for every distinct pair. Three
distinct labels give $d\mid30a_i$ for every $i$. Since $d$ is odd and
the word is primitive, this implies $d\mid15$, a contradiction.
Choose a pair $i,j$ whose eigenvalue has order greater than $5$.

If some third label gives a nonzero cross-sum, it gives the desired
triple, since every nontrivial root of odd order has order greater than
$2$. Otherwise $a_i=a_j=a$ and every other entry equals $-a$.
For $s\ge5$, take three of those other entries: their pair sums are
$-2a$, of the same order as $2a=a_i+a_j$, so they give a good triple.
For $s=4$, the word is $(a,a,-a,-a)$; its age equals $2$ at every
primitive embedding and it is mixed.

Orders and the vanishing of cross-sums are unchanged under multiplying
the word by a unit. At a mixed embedding $2\le Q_a(b)\le s-2$.
Choose an active label as infinity. Its omission leaves $n=s-1$ finite
weights whose sum is strictly between $1$ and $n-1$, as required by
\cite[Definition~1.1(a)]{MN}. At a definite embedding the age is
$1$ or $s-1$; choose a good triple and a different label as infinity.
The finite weight sum is respectively less than $1$ or greater than
$n-1$, and the triple meets \cite[Definition~1.1(b)]{MN}.
Primitivity persists after omitting the infinity entry, because that
entry is the negative sum of the finite ones. Theorem~5.1 of
\cite{MN} therefore applies. Its real connected group is the full
special unitary group; its complexification is $\mathrm{SL}_h$.
\end{proof}

For $d=3,5$, the prime-conductor argument of \cite[Section~2]{General}
applies: every word
with $s\ge4$ has a mixed primitive embedding, and Galois conjugation
transports its full complex monodromy to all primitive embeddings.
The same Galois argument will be used below: in a cyclotomic basis the
monodromy matrices have entries in $\mathbb Q(\zeta_d)$, and the
Zariski-closure ideal commutes with scalar extension. Applying a field
automorphism transports the full $\mathrm{SL}_h$ closure.

\begin{lemma}[The unique all-definite conductor-fifteen word]
\label{lem:fifteen-exception}
For $d=15$ and $s\ge4$, all primitive projections have full
$\mathrm{SL}_{s-2}$ monodromy unless, up to a unit multiple and a
permutation, $a=(1,2,4,8)$. This exceptional primitive rational factor
is isogenous to $E^8$, where $E$ is a fixed elliptic curve with complex
multiplication by $\mathbb Q(\sqrt{-15})$. Moreover $E$ is an isogeny
factor of the degree-fifteen Fermat Jacobian.
\end{lemma}
\begin{proof}
If some pair eigenvalue has order greater than $5$, the preceding
good-triple argument applies, including its mixed four-point case.
It remains to consider words whose pair sums are all divisible by $3$
or $5$.

Such a primitive word contains a unit: otherwise it must contain both
a nonzero multiple of $3$ and a nonzero multiple of $5$, and their sum
would be a unit. Normalize one unit entry to $1$. The other entries
then belong to
\[
 \{2,4,5,8,9,11,14\}.
\]
A unit cannot be repeated, since twice a unit is a unit. Directly
testing divisibility by $3$ or $5$ shows that the maximal compatible
support sets containing $1$ are
\[
 \{1,2,4,8\},\quad \{1,4,11,14\},\quad
 \{1,4,5\},\quad \{1,9,11\}.
\]
For example, $5$ is compatible only with the units $1,4$ in this list,
and $9$ only with $1,11$; they are incompatible with each other.
Among the units, the first two displayed sets are the maximal cliques.
This is a seven-vertex finite divisibility check, not a numerical
search over branch configurations.

If all entries are units, $s\ge4$ forces one of the two four-element
sets. The word $(1,4,11,14)$ is two opposite pairs and is mixed.
For support involving $5$, both units $1,4$ must occur: with only one
unit the sum cannot vanish modulo $5$. The word therefore is
$(1,4,5^m)$, with $m=2+3t$ and $t\ge0$. At embedding $b=2$, its Hodge
multiplicities are $(1+t,1+2t)$, both positive.
For support involving $9$, the word similarly is $(1,11,9^m)$,
$m=2+5t$; at embedding $b=1$ its multiplicities are
$(1+2t,1+3t)$. Galois propagation now gives full complex monodromy in
all these mixed cases.

For the remaining word $(1,2,4,8)$, put
$K=\mathbb Q(\zeta_{15})$ and $P=\langle2\rangle=\{1,2,4,8\}$.
The primitive eigenspaces have multiplicities $(2,0)$ on $P$ and
$(0,2)$ on $-P$. Write
\[
 \theta=\zeta_{15}+\zeta_{15}^2+\zeta_{15}^4+\zeta_{15}^8.
\]
The sum over all primitive fifteenth roots gives
$\theta+\bar\theta=1$. The twelve ordered differences of distinct
elements of $P$ are exactly the nonzero residues other than $5,10$;
hence
\[
 \theta\bar\theta
 =4+\sum_{j\ne0,5,10}\zeta_{15}^j=4.
\]
Thus $\theta^2-\theta+4=0$, and
$K^P=\mathbb Q(\sqrt{-15})=:K_0$.

Choose an elliptic curve $E$ with CM by $K_0$ and the type whose
extensions to $K$ form $P$. The weight-one rational Hodge structure
\[
 T=H^1(E,\mathbb Q)\otimes_{K_0}K
\]
has rank one over $K$ and exactly that pure type at each embedding.
The exceptional primitive factor has rank two over $K$, so any
$K$-linear identification with $T^{\oplus2}$ is a Hodge isomorphism:
each embedding summand has just one weight-one Hodge type.
Since $[K:K_0]=4$, its rational first cohomology is
$H^1(E,\mathbb Q)^{\oplus8}$. The full faithfulness of rational
weight-one Hodge structures for complex abelian varieties up to
isogeny gives an actual isogeny with $E^8$.
In particular this supplies algebraic maps on the entire primitive
first cohomology, not merely a numerical determinant-type match.
Multiplying the residue word by a unit can exchange the two $K_0$
types; conjugation on $K_0$ gives a rational Hodge isomorphism between
these realizations after forgetting the specified $K_0$ action, so
the same elliptic isogeny class serves both types.
At a fixed fiber these are the algebraic maps required
by the invariant-theory assembly; no global choice of an isogeny is
needed. The primitive rational summand has a monodromy-stable integral
lattice, obtained by intersecting it with integral cohomology. Every
real factor of its polarization-preserving cyclotomic centralizer is
compact, because every primitive embedding is definite. The monodromy
therefore lies in the intersection of a compact real group with the
discrete group of integral lattice automorphisms and is finite.

Finally, let $C_0$ normalize $y^{15}=x(x-1)^2$, whose residues are
$(1,2,12)$. Its primitive rank-one $K$-factor has multiplicities
$(1,0)$ on $P$ and $(0,1)$ on $-P$. Thus its primitive abelian factor
is isogenous to $E^4$. On the Fermat curve $u^{15}+v^{15}=1$, the map
\[
 x=u^{15},\qquad y=uv^2
\]
gives a finite covering of $C_0$. Pullback and norm therefore make
$E$ an isogeny factor of the Fermat Jacobian.
All statements about $E^8$ and $E^4$ concern the primitive factors:
the whole four-branch and three-branch Jacobians have dimensions
$14$ and $6$, respectively.
\end{proof}

\subsection{The exact simultaneous blocks at varying conductors}
Write $M$ for the connected algebraic monodromy group on the full
rational first cohomology. The projection calculations above were made
on affine configuration spaces. Normalizing two finite points identifies
$\operatorname{Conf}_{s-1}(\C)$ with
$\operatorname{Aff}_1\times M_{0,s}$. An affine coordinate change
rescales the Kummer variable by a root of a power of its scale factor,
so the affine fiber acts by finite deck scalars on every character.
Thus each projective character representation descends to $M_{0,s}$
with the full connected projection unchanged. Forgetting inactive
labels, $M_{0,r}\to M_{0,s}$ has connected configuration-space fibers
and surjects on fundamental groups. Support-three primitive factors
are constant after a finite root choice, by normalizing their three
active points to $0,1,\infty$. Pair-braid determinants and the scalar
ambiguities have finite order by \cite[Corollary~2.7]{MN}.

For the indices with $s(c)\geq4$ other than the finite factors of
Lemma~\ref{lem:fifteen-exception}, the complex blocks are exactly those of
\cite[Section 3]{General}, with words now read in $\Z/N$:
different supports are independent; on support at least five the block
is $\{c,-c\}$; on support four it is
\[
 \bigl(\mathcal V_4(S)c\cup-\mathcal V_4(S)c\bigr)\cap C.
\]
We check the changes needed for this extension.

The connected solvable radical of $M_\C$ maps to a connected solvable
normal subgroup of each full $\SL_h$ projection, hence trivially.
It also acts trivially on the finite-monodromy factors. Faithfulness
therefore makes the radical trivial, and semisimple Goursat theory
reduces the product determination to the pairwise graph relations.

All projective character representations are defined on the same
$\pi_1(M_{0,r})$, before a finite base cover: two lifts of a marked
sphere mapping class differ by a deck transformation, scalar on each
character. This asserts a product of projective representations, not
one common scalar on their direct sum. In a product of two
centerless simple groups the normalizer of a graph of an automorphism
is that graph. Hence a graph relation between connected projections
constrains the original pair braids, before any powers kill their finite
eigenvalues. For an active pair, the rank-one formula of
\cite[Theorems 2.2 and 2.5]{MN}, with a different active point at
infinity, is nonzero: the $n$ spanning vectors have their sole relation
$\sum g_i=0$, so its vector is nonzero, and the Hermitian form and its
coefficient are nondegenerate and nonzero. The projective twist is
therefore nonidentity even when the pair sum is zero, in which case it
is unipotent. If a label is inactive for the other character, forgetting
it kills the other character's pair twist. This proves support separation.

On a fixed support, the exceptional pair eigenvalue is
$\zeta_N^{c_i+c_j}$, equivalently
$\zeta_{d(c)}^{b_i+b_j}$. In rank $h\geq3$, its multiplicities
$1,h-1$ and the inner-or-dual form of an automorphism of $\PGL_h$
give one global sign $\epsilon$ and all equations
$e_i+e_j=\epsilon(c_i+c_j)$ modulo $N$. Since $2$ is invertible
modulo $N$, three finite indices reconstruct $e_i=\epsilon c_i$;
the sum-zero equation reconstructs the infinity coordinate.
In rank two, put $x=c_1+c_2$, $y=c_1+c_3$, $z=c_1+c_4$. The inverse
\[
 c=\tfrac12(x+y+z,\ x-y-z,\ -x+y-z,\ -x-y+z)
 \quad\text{in }(\Z/N)^4
\]
identifies the eight independent sign changes of the three unordered
pair ratios with $\{\pm P:P\in\mathcal V_4\}$. Zero sums cause
only stabilizers. These are necessity arguments over the common ring
$\Z/N$; no field assumption is hidden in the reconstruction.

Conversely $e=Pc$ preserves the additive order $d$ of the word.
Normalize the four branch points to $0,1,t,\infty$; the three
double transpositions are realized by
\[
 T_{12|34}(x)=\frac{t(x-1)}{x-t},\qquad
 T_{13|24}(x)=\frac{x-t}{x-1},\qquad T_{14|23}(x)=\frac{t}{x}.
\]
Dividing the two words by $N/d$ gives degree-$d$ Kummer equations
$y^d=f_c(x)$ and $z^d=f_e(x)$. The divisor of
$f_e(T_P(x))/f_c(x)$ is divisible by $d$; a degree-zero divisor on
$\mathbf P^1$ is principal. Consequently the quotient is
$a(t)h(x)^d$, where $a(t)$ is a constant times powers of $t$ and
$t-1$. Adjoining $a(t)^{1/d}$ gives the deck-equivariant isomorphism
$(x,y)\mapsto(T_P(x),a(t)^{1/d}h(x)y)$ on the smooth normalizations.
The finite root cover is etale on $M_{0,4}$. Pullback and
norm, followed by \eqref{eq:ppprimitive}, give algebraic isomorphisms
of their primitive factors at every primitive embedding. Thus
$O_c=(\mathcal V_4(S)c)\cap C$ is the actual oriented component;
$-O_c$ is its polarized dual, possibly equal to $O_c$. The negative
orientation uses polarization as in \cite{General}; it is not silently
replaced by a degree-zero Hodge isomorphism. Thus no graph can identify
different conductors. Goursat and the faithful standard occurrences give
the stated product of $\SL_h$ factors, with no additional connected
scalar factor or hidden pairwise relations.
Multiplication of all character words by a unit modulo $N$ preserves
the description, so it is Galois-stable.

\subsection{Primitive determinants as Fermat Chow summands}
Let $C_c$ have odd degree $d$, active support $s=h+2\geq3$, and primitive
factor $U=U_{[c]}$. Its rational determinant is
\[
 D_c\otimes\C=
 \bigoplus_{u\in(\Z/d)^\times}\bigwedge^h V_{uc},
 \qquad \dim_\Q D_c=\varphi(d).
\]
It is realized in both $H^h(C_c^h,\Q)$ and $H^h(J(C_c),\Q)$.
The finite-quotient determinant construction at arbitrary degree,
including the primitive cyclotomic factor and the deck action in one
slot, already appears in Schoen \cite[Lemmas~1.1--1.2, pp.~5--6]{Schoen}.
We record its primitive projector and the normalization of the two
correspondences explicitly, extending the presentation in \cite{Fermat};
the composite-degree mechanism is not claimed as a new construction.

Write $y^d=\prod_{i=1}^{h+2}(x-t_i)^{b_i}$ with all active points
finite and $\sum b_i=0$ modulo $d$. Put $Z=C_c^h$ and
\[
 H=\ker(\mu_d^h\xrightarrow{\prod}\mu_d)\rtimes\mathfrak S_h,
 \qquad \widetilde H=\mu_d^h\rtimes\mathfrak S_h.
\]
The quotient $Y=Z/H$ has function field
\[
 \C(\Sym^h\mathbf P^1)(z),\qquad
 z^d=\prod_i L_i^{b_i},
\]
where $L_i$ is evaluation at $t_i$. The Kummer order is exactly $d$
because $\gcd(d,b_1,\ldots,b_s)=1$. Likewise the independent coordinate
Kummer classes on the ordered product have order $d$, as their
coordinate valuations show. The $h+2$ evaluation hyperplanes are in
general position, with a unique relation having all coefficients nonzero.
Taking $d$th roots of these coordinates produces a smooth diagonal
Fermat hypersurface $F_d^h$ finite over $\mathbf P^h$.
The character $(\epsilon_i)\mapsto\prod_i\epsilon_i^{b_i}$ of
$\mu_d^{h+2}/\mu_d$ is surjective, with kernel of order $d^h$.
Its quotient has the same function field as $Y$; normality identifies
the quotients globally. Thus there are finite maps
\[
 F_d^h\xrightarrow{r}Y\xleftarrow{q}Z,
 \qquad d_r=d^h,\quad d_q=d^{h-1}h!.
\]

Let $e_H,e_{\widetilde H}$ be the rational averaging correspondences.
The group $H$ is the kernel of the product map
$\widetilde H\to\mu_d$, hence is normal. A deck generator $T$
acting in one coordinate therefore commutes with both averages.
The following is a rational Chow idempotent on $Z$:
\begin{equation}\label{eq:ppdetprojector}
 P_d=e_d(T)(e_H-e_{\widetilde H}).
\end{equation}
On cohomology, the kernel-of-product average forces all deck labels
to agree. The primitive projector forces that label to be primitive
and nonzero, so every factor is in $H^1(C_c)$; $H^0$ and $H^2$
have trivial deck label. The unsigned geometric permutation average
acts on these degree-one tensors as the ordinary antisymmetrizer,
by the Koszul sign. Consequently the image of $P_d$ is concentrated
in degree $h$ and is exactly $D_c$. Lower-conductor eigenspaces of
the entire cyclic quotient are removed by $e_d(T)$.

For $\mathcal T=[F_d^h\times_Y Z]$, the finite-correspondence identity is
\[
 \mathcal T\mathcal T^{\mathsf t}
   =d_r\sum_{g\in H}\Gamma_g.
\]
This holds in rational Chow groups. The relevant intersection is proper
and finite over the dimension-$h$ base; its generic components are the
graphs in the displayed sum. Components supported over the branch locus
have dimension less than $h$ and cannot contribute to that cycle.
It follows that
\[
 R_d=(d_r d_q)^{-1}\mathcal T^{\mathsf t}P_d\mathcal T
\]
is a rational algebraic idempotent on $F_d^h$. The maps
$P_d\mathcal T$ and $(d_r d_q)^{-1}\mathcal T^{\mathsf t}P_d$
identify its image with the $P_d$ summand in both directions.
The Jacobian realization uses the algebraic degree-one projector,
the Abel--Jacobi $H^1$ isomorphism and its inverse, and alternation.
Choose any basepoint on $C_c$. A composite cyclic cover need not have
a point fixed by its whole deck group, but changing the basepoint
changes Abel--Jacobi by a translation, which acts trivially on $H^1$.
Thus its $H^1$ isomorphism intertwines the deck actions without a
fixed-point assumption, as observed in
\cite[Remark~1.4, pp.~6--7]{Schoen}. The inverse weight-one Hodge map is algebraic
by the divisor realization of homomorphisms of Jacobians. The addition
and diagonal formulas give the same nonzero $h!$ normalization as in
\cite{Fermat}; none depends on $d$ being prime. The asserted Chow
idempotent identities remain on $C_c^h$ and $F_d^h$; the Jacobian
comparison here is an algebraic isomorphism of the specified
cohomological summands.

\subsection{The precise Fermat premise and its established cases}
Assume the rational Hodge conjecture on every power of $J(F_N^1)$.
We need arbitrary products of Fermat hypersurfaces, and record the
reduction instead of silently treating a theorem about curves as a
theorem with that different statement. Aoki's original
\cite[Theorem 0.1, p.~177]{Aoki2000} supplies this premise when $N$
is a power of an odd prime, or when $N=3^a5^b$ or $N=3^a7^b$
with $a,b\geq0$ and $N>1$: his assertion covers all abelian varieties
isogenous to factors of powers of the corresponding Fermat Jacobian.

First the conjecture on $J(F_N^1)^q$ implies it on $(F_N^1)^q$:
Abel--Jacobi pullback surjects on the cohomology of the curve in degrees
$0,1,2$, hence on the product, and semisimplicity lifts rational Hodge
classes through this surjection. Pullback of the resulting algebraic
classes is algebraic.

Write $F_N^a=\{\sum_{i=0}^{a+1}x_i^N=0\}$. For $a,b\geq1$,
the Shioda--Katsura construction \cite[Lemma 1.2 and Theorem 1.7]{SK}
resolves the dominant degree-$N$ rational map
\[
 F_N^a\times F_N^b\dashrightarrow F_N^{a+b},\qquad
 (x,y)\longmapsto
 [x_0y_{b+1}:\cdots:x_a y_{b+1}:
   \eta x_{a+1}y_0:\cdots:\eta x_{a+1}y_b],
 \quad\eta^N=-1,
\]
by blowing up the smooth codimension-two center
$F_N^{a-1}\times F_N^{b-1}$ defined by
$x_{a+1}=y_{b+1}=0$. Indeed the base ideal is locally generated by
these two coordinates; the blowup removes their common factor in the
displayed sections. Their $N$th powers sum to zero by the defining
equations. Away from the center, recovering the two original factors
leaves exactly $N$ choices. Thus the resolved morphism $f$ satisfies
$f_*f^*=N\id$ on rational cohomology.

Apply this with $a=h-1,b=1$ to a Fermat factor of dimension $h\geq2$,
and take its product with all the other factors. The blowup formula
expresses its cohomology algebraically as that of the product with
$F_N^{h-1}\times F_N^1$, plus the center
$F_N^{h-2}\times F_N^0$ with a Tate twist. Induction on
$\sum_j\max(h_j-1,0)$, starting with curve products and points,
proves the Hodge conjecture for every product $\prod_jF_N^{h_j}$.
Here $F_N^0$ is a disjoint union of $N$ points. At every stage both
summand maps are algebraic, and $N^{-1}f_*$ descends algebraic lifts.
Finally, if $d\mid N$, the finite power map
\[
 F_N^h\longrightarrow F_d^h,\qquad [x_i]\longmapsto[x_i^{N/d}]
\]
gives the same conclusion for arbitrary products with degrees dividing
$N$. The primitive determinant transfers just proved imply that every
rational Hodge class in every full rational tensor product of the
$D_c$ is algebraic in its specified realization. In addition, if $E$
is the constant elliptic factor supplied by the conductor-fifteen
classification, then $H^1(E,\Q)$ is an algebraic direct summand of
$H^1(F_{15}^1,\Q)$, and hence of $H^1(F_N^1,\Q)$ whenever $15\mid N$.
An isogeny factor of the Fermat Jacobian gives the two $H^1$
correspondences by the isogeny projection and its rational inverse,
composed with Abel--Jacobi; equality of Hodge numbers alone would not
suffice. Therefore the same conclusion holds for mixed tensors of the
$D_c$ and any number of copies of $H^1(E,\Q)$, after effective Tate
twists and rational algebraic projections as well.

\subsection{Rational invariant assembly with different conductors}
It remains to explain why mixed conductors require no unsupported
matching-embedding assertion. Use the full rational sources
\[
 E_{c_1,\ldots,c_t;\,u}=
 D_{c_1}\otimes_\Q\cdots\otimes_\Q D_{c_t}
                  \otimes_\Q H^1(E,\Q)^{\otimes u},
\]
including repeated entries and the empty source $\Q$. Here $u=0$
unless the exceptional conductor-fifteen factors occur; each such
primitive factor is supplied by the classification as an actual
weight-one Hodge structure isomorphic to $H^1(E,\Q)^{\oplus8}$.
The corresponding abelian varieties are isogenous because rational
weight-one Hodge morphisms between complex abelian varieties are
induced by homomorphisms up to isogeny. Thus these exceptional slots
are inserted by rational algebraic maps.
Such a source generally is \emph{not} rank one over $\Q(\zeta_N)$.
That property is neither claimed nor needed.

For each of the product-$\SL_h$ blocks, tensor invariant theory expresses
all invariants through pair contractions and determinant volumes.
Actual oriented graphs transfer occurrences to representatives;
polarization handles their dual occurrences. This includes the
self-dual rank-two blocks and the support-three rank-one factors.
The exceptional rank-two primitive factors are constant after finite
base change, so their slots are unrestricted tensors of the additional
$H^1(E)$ source. A diagram of total degree $2k$ therefore gives an
algebraic insertion map from some $E_{c_1,\ldots,c_t;u}(-v)$, with
$\sum h(c_i)+u+2v=2k$ and $v\geq0$, into $H^{2k}(A^q,\Q)$.

For precision about rational spanning, work temporarily over
$L=\Q(\zeta_N)$. Every desired primitive character selector on a raw
$H^1$ slot is the Fourier idempotent
\[
 e_c=\frac1{|G|}\sum_{g\in G}\zeta_N^{-c(g)}g^*
 \quad\text{in }L[G].
\]
It is an $L$-linear combination of rational algebraic deck
correspondences. Each required determinant line is a component of its
full rational $D_c$. A contraction is obtained by applying the two
character selectors to the rational polarization tensor, and an
oriented identification by applying selectors to its actual rational
graph correspondence. Alternation and passage to exterior cohomology
are rational algebraic, including their factorial normalizations.
Thus, on expanding these finitely many Fourier selectors, every
$L$-valued invariant diagram belongs to the $L$-linear span of images
of rational algebraic maps from the full sources $E(-v)$ above.
All these maps intertwine connected monodromy; the determinant sources
have finite determinant monodromy and the additional elliptic sources
map to constant primitive factors. Their images therefore lie in the
invariant part.
There are only finitely many diagrams, deck terms and source choices
in a fixed cohomological degree. The resulting finite sum is a rational
Hodge surjection
\[
 \bigoplus_\alpha E_\alpha(-v_\alpha)
    \twoheadrightarrow H^{2k}(A^q,\Q)^M.
\]
Surjectivity follows after extension to $L$ from the spanning just
proved, hence holds over $\Q$. This argument also supplies the
coefficients needed to distinguish conjugate contraction slots; it does
not merely average one chosen diagram or assume a rational monomial
basis. Equivalently, one may refine the sources by the rational
Galois-orbit idempotents in their tensor products of cyclotomic deck
algebras, but that refinement is unnecessary.

For a very general fiber, connected monodromy is contained in its
Hodge group, so every rational Hodge class lies in the displayed target.
Polarizable rational Hodge structures are semisimple; a Hodge class
therefore lifts to a rational Hodge class in the source. The preceding
Fermat argument makes every such source class algebraic, and the
insertion maps preserve algebraicity. This proves the required
surjectivity of the cycle class map in every degree and on every power.

Finally the very-general assertion uses exactly the monodromy--Hodge
normality and algebraic Hodge-locus argument in \cite[Section 6]{General},
which has no prime-exponent hypothesis. Apply the Hodge-locus theorem
after the appropriate weight-zero Tate twist and take the countable
union over powers and degrees. Images of these proper closed loci under
the finite base map remain proper and closed. The conclusion descends
to the original full branch family. Together with the simultaneous-block
and exceptional-constant classification proved above, this completes
Theorem~\ref{thm:odd-exponent-criterion}.

\begin{corollary}[Aoki exponents]
\label{cor:aoki-odd-exponents}
The conclusion of Theorem~\ref{thm:odd-exponent-criterion} holds
unconditionally on its stated very general locus if
\[
 N=p^e\quad(p\text{ odd prime},\ e\ge1),\qquad
 N=3^a5^b,\qquad\text{or}\qquad N=3^a7^b,
\]
where $a,b\ge0$ and $N>1$.
\end{corollary}
\begin{proof}
Aoki's \cite[Theorem~0.1(ii)]{Aoki2000} applies to every odd prime power.
Theorem~0.1(i), for
$m=2^a3^b5^c7^d$ with $c=0$ or $d=0$, gives the other two odd
families. His definition of Fermat type includes every isogeny
factor of every power of the Fermat Jacobian, so it supplies
precisely the hypothesis of the criterion.
\end{proof}


\section{Verification and prior-art boundary}
The exceptional-space arithmetic was independently checked for all odd
subgroups of prime residue groups with $p\le103$, with zero through three
added opposite pairs and $s\ge5$: 152 cases. The checks verify branch
balance, signature rows, generator inversion, diagonal injectivity and
inverse-parameter duplicates. A further 7,746 volume words of lengths two
and four at $p=7,13,19$ verify the opposite-coset balance criterion.
The odd-exponent regressions check eleven exact rational primitive
projectors, six finite-abelian character inventories, and 24,159
primitive support-four words at odd conductors at most 45, together
with the conductor-fifteen CM type and a cover without a deck-fixed
basepoint. Normal and optimized Python give identical reports.
They do not verify algebraic cycles,
deformation, monodromy, or the universal quantifiers proved above.

The determinant-to-Fermat construction has direct precedent in Schoen's
1989 work, as documented by \cite{Fermat}; the algebraicity input there is
Shioda's theorem on products of prime-degree Fermat varieties. The alternative
simple-tuple cycle construction is Schoen's
\cite{Schoen,SchoenAddendum}, extended by the standalone fusion proof.
Although the moving Jacobians in Section~\ref{sec:exceptional} are
non-CM, their determinant motives are direct summands of Fermat motives.
Their exceptional-cycle algebraicity is therefore a consequence of this
classical construction and Shioda, not a new algebraicity theorem obtained
by observing that the entire Jacobian is non-CM. The three-point boundary
is itself a Fermat quotient and is excluded from the moving simple assertion;
see also the prime-Fermat results summarized in Aoki \cite{AokiCM}.
Patel--Zhang \cite{PZ} extend Schoen's construction for
\emph{etale} abelian covers and explicitly leave ramified cases to future
work; an exact comparison of alternative realizations remains necessary.
Rohde \cite[Theorem 5.3.7]{Rohde} already gives the full prime-cyclic
monodromy input, making the single-cyclic ordinary all-powers consequence
a deduction from published group and cycle results. Spelta--Tamborini
\cite[Theorem 1.7 and Section 3]{ST} treat simultaneous abelian-cover
monodromy under a nonrepetition condition and analyze examples with
repeated factors; their results are direct precedent for the independence
question. The signature-orbit mechanism in Section~\ref{sec:hg} is
standard and appears explicitly in Xue--Zarhin \cite{XZ}.
The degree-seven generalized-Hodge corollary is attributed in
Section~\ref{sec:ghc7}, including its extra CM and domination steps.
The odd-exponent extension in Section~\ref{sec:odd-exponent} uses
Menet--Nguyen's good-weight criterion and an explicit treatment of the
constant conductor-fifteen exception. Its cycle transfer uses the
classical Shioda--Katsura construction and Aoki's all-powers Fermat
theorem. Schoen's \cite[Lemmas~1.1--1.2 and Remark~1.4]{Schoen}
already supplies the arbitrary-degree primitive determinant mechanism
and the Abel--Jacobi translation observation. The small residue
computations are not premises of this
extension, and the extension makes no all-odd-exponent claim without
its stated Fermat premise.
None of these comparisons
constitutes a priority claim for the statements of this paper.

All results depend on the companion geometric proofs and their published
inputs. The existing Lean files formalize only designated arithmetic and
abstract interfaces. Independent specialist review of the mathematical
manuscripts remains necessary. The generalized assertion has the additional
published domination input and the explicit quadratic CM verification;
ordinary all-powers algebraicity alone is not substituted for them.

\begin{thebibliography}{99}
\bibitem{General} K. Kistner, \emph{All powers of very general Jacobians of
elementary abelian covers}, 5 September 2026. Companion source
\nolinkurl{manuscripts/general_all_powers.tex}.
\bibitem{Fusion} K. Kistner, \emph{Balanced mixed determinants of prime
cyclic covers: a revised fusion argument}, 4 September 2026. Companion
source \texttt{manuscripts/fusion\_revised.tex}, Theorem 1.1 and Sections 2--6.
\bibitem{Fermat} K. Kistner, \emph{Determinants of cyclic covers of the line:
a finite-quotient reduction to Fermat varieties}, 5 September 2026.
Companion source \nolinkurl{manuscripts/fermat_determinant_transfer.tex},
Theorem 1.1 and Corollary 1.2.
\bibitem{MN} G. Menet and D.-M. Nguyen, \emph{Representations of braid groups
via cyclic covers of the sphere: Zariski closure and arithmeticity},
arXiv:2310.10401v3, \url{https://arxiv.org/html/2310.10401v3}.
\bibitem{XZ} J. Xue and Y. G. Zarhin, \emph{Centers of Hodge groups of
superelliptic Jacobians}, preprint, version 4,
\href{https://arxiv.org/pdf/0907.1563}{arXiv:0907.1563},
Lemma 3.9, Theorem 3.11 and Section 3.13.
\bibitem{Schoen} C. Schoen, \emph{Hodge classes on self-products of a variety
with an automorphism}, Compositio Math. 65 (1988), 3--32,
\url{https://www.numdam.org/item/CM_1988__65_1_3_0.pdf}.
\bibitem{SchoenAddendum} C. Schoen, \emph{Addendum to: Hodge classes on
self-products of a variety with an automorphism}, Compositio Math. 114
(1998), 329--336, \url{https://doi.org/10.1023/A:1000566205021}.
\bibitem{AokiCM} N. Aoki, \emph{Hodge cycles on CM abelian varieties of
Fermat type}, Comment. Math. Univ. St. Pauli 51 (2002),
\url{https://rikkyo.repo.nii.ac.jp/record/8744/files/AA00610867_51-01_07.pdf}.
\bibitem{PZ} D. Patel and Y. Zhang, \emph{Algebraicity of Hodge classes on
some generalized Prym varieties}, arXiv:2506.13729v2,
\url{https://arxiv.org/html/2506.13729v2}.
\bibitem{Rohde} J. C. Rohde, \emph{Cyclic coverings, Calabi--Yau manifolds
and Complex multiplication}, dissertation, Universit\"at Duisburg--Essen,
2007, Construction 3.2.1 and Theorem 5.3.7 (printed p.\ 64),
\url{https://webdoc.sub.gwdg.de/ebook/dissts/Duisburg/Rohde2007.pdf}.
\bibitem{ST} I. Spelta and C. Tamborini, \emph{Decomposable abelian
$G$-curves and special subvarieties}, Bull. Sci. Math. 201 (2025), 103616,
\href{https://doi.org/10.1016/j.bulsci.2025.103616}{doi:10.1016/j.bulsci.2025.103616};
\url{https://arxiv.org/pdf/2305.19030v2}.
\bibitem{Shioda79} T. Shioda, \emph{The Hodge conjecture and the Tate
conjecture for Fermat varieties}, Proc. Japan Acad. Ser. A Math. Sci. 55
(1979), 111--114, Theorem 2, p.\ 112,
\href{https://doi.org/10.3792/pjaa.55.111}{doi:10.3792/pjaa.55.111}.
\bibitem{Abdulali12} S. Abdulali, \emph{Tate twists of Hodge structures
arising from abelian varieties of type IV}, J. Pure Appl. Algebra 216
(2012), 1164--1170, Proposition 1, Lemma 9, Theorems 10--11 and Remark 12,
\href{https://doi.org/10.1016/j.jpaa.2011.10.005}{doi:10.1016/j.jpaa.2011.10.005};
\url{https://arxiv.org/html/1203.4857v1}.
\bibitem{AbdulaliSurvey} S. Abdulali, \emph{Tate twists of Hodge structures
arising from abelian varieties}, in Recent Advances in Hodge Theory:
Period Domains, Algebraic Cycles, and Arithmetic, London Math. Soc.
Lecture Note Ser. 427, Cambridge Univ. Press, 2016, 292--307,
Appendix B.5, \url{https://myweb.ecu.edu/abdulalis/Vancouver.pdf}.
\bibitem{Aoki2000} N. Aoki, \emph{Some remarks on the Hodge conjecture for
abelian varieties of Fermat type}, Commentarii Mathematici Universitatis
Sancti Pauli \textbf{49} (2000), 177--194, especially Theorem 0.1(i)--(ii),
p.~177. \url{https://rikkyo.repo.nii.ac.jp/record/9851/files/AA00610867_49-02_04.pdf}.
\bibitem{SK} T. Shioda and T. Katsura, \emph{On Fermat varieties},
Tohoku Mathematical Journal \textbf{31} (1979), 97--115, especially
Lemma 1.2 and Theorem 1.7.
\url{https://www.jstage.jst.go.jp/article/tmj1949/31/1/31_1_97/_pdf/-char/en}.
\end{thebibliography}
\end{document}
